\chapter{Solution of Linear Algebraic Equations}
\label{ch:solve}

\section{Methods for solving linear algebraic equations \texorpdfstring{$\mat{A}\mat{x}=\mat{b}$}{Ax=b}}
\begin{enumerate}
  \item Gauss elimination
  \item Decomposition methods
  \item Iterative methods: such as Gauss--Seidel and Jacobi methods and conjugate gradient method (Sec.~11.3 of textbook).
\end{enumerate}
We shall discuss only the first two methods because they are the most commonly used methods.

\section{Gauss elimination method}
The main idea of Gauss elimination is to reduce the original system $\mat{A}\mat{x} = \mat{b}$ to an upper triangular system through appropriate \emph{row operations}; that is, the desired result is
\begin{equation*}
  \mat{U}\mat{x} = \mat{c}
\end{equation*}
where $\mat{U}$ is an upper triangular matrix. This system can then be solved to obtain $\mat{x}$ using the backward substitution. Before explaining the algorithm, let's consider the following simple example.

\begin{example}
Solving the linear system with Gauss elimination.
\begin{alignat*}{2}
  3x_1 - 6x_2 + 7x_3 &= 3 &\qquad&\text{(a)}\\
  9x_1 + 0 - 5x_3 &= 3 &&\text{(b)}\\
  5x_1 - 8x_2 + 6x_3 &= -4 &&\text{(c)}
\end{alignat*}
First, from (a): $x_1 = (6x_2-7x_3+3)/3 \;\rightarrow\;$ (b), (c). Then, we have
\begin{alignat*}{2}
  18x_2 - 26x_3 &= -6 &\qquad&\text{(b$'$)}\\
  2x_2 - \tfrac{17}{3}x_3 &= -9 &&\text{(c$'$)}
\end{alignat*}
Second, from (b$'$): $x_2 = \frac19(-3+13x_3) \;\rightarrow\;$ (c$'$). Then, we have
\begin{alignat*}{2}
  -\tfrac{25}{9}x_3 &= -\tfrac{25}{3} &\qquad&\text{(c$''$)}
\end{alignat*}
Rewrite these results [(a), (b$'$), and (c$''$)] in matrix form:
\begin{equation*}
  \bmat{3 & -6 & 7\\ 0 & 18 & -26\\ 0 & 0 & -\frac{25}{9}}\Bmat{x_1\\ x_2\\ x_3} = \Bmat{3\\ -6\\ -\frac{25}{3}}
\end{equation*}
The above procedure is called \emph{forward elimination}, which results in an upper triangular system $\mat{U}\mat{x} = \mat{c}$.

The upper triangular system can be solved easily using the \emph{backward substitution} as described below. First, obtain $x_3 = 3$ from (c$''$). Then substituting $x_3 = 3$ into (b$'$) gives $x_2 = 4$. Finally, substitute $x_3 = 3$ and $x_2 = 4$ into (a) to obtain $x_1 = 2$.
\end{example}

\topic{Algorithm for upper triangularization: $\mat{A}\mat{x}=\mat{b}\ \rightarrow\ \mat{U}\mat{x}=\mat{c}$.}
\begin{lstlisting}[language=Fortran,numbers=none,basicstyle=\small\ttfamily\itshape,mathescape=true]
do m=1, n-1
    do i=m+1, n
        r=$a_{im}/a_{mm}$
        $b_i$=$b_i$-r$b_m$
        do j=m, n
            $a_{ij}$=$a_{ij}$ - r $a_{mj}$
        end do
    end do
end do
\end{lstlisting}
It can be represented by the following graph.
\begin{fig}
  % Graphical representation of the upper triangularization step m
\begin{tikzpicture}[scale=0.85,font=\small]
  \def\n{8}
  % matrix frame
  \draw (0,0) rectangle (\n,\n);
  % processed upper-left triangle: X on the diagonal (rows 1..m-1)
  \foreach \i in {1,...,5}{\node at (\i-0.5,\n-\i+0.5) {X};}
  \draw[dashed] (0.6,\n-0.6) -- (2.7,\n-2.7); \node at (3.0,\n-3.0) {$\cdot$}; \node at (3.3,\n-3.3) {$\cdot$};
  % dashed region of zeros below the processed part
  \draw[dashed] (0.3,\n-0.3) -- (0.3,0.3) -- (3.7,0.3) -- (3.7,\n-3.7);
  \node[font=\Large] at (1.6,3.6) {$0$};
  % pivot row m and column m, column j, row i
  \pgfmathsetmacro{\m}{\n-5.5} \pgfmathsetmacro{\ii}{\n-7.0}
  \draw[dotted] (0,\m) -- (\n,\m); \node[left] at (0,\m) {$m$};
  \draw[dotted] (0,\m-0.35) -- (3.7,\m-0.35); \node[left] at (0,\m-0.35) {$m{+}1$};
  \draw[dotted] (0,\ii) -- (\n,\ii); \node[left] at (0,\ii) {$i$};
  \draw[dotted] (3.7,0) -- (3.7,\n); \node[above] at (3.7,\n) {$m$};
  \draw[dotted] (6.3,0) -- (6.3,\n); \node[above] at (6.3,\n) {$j$};
  \node[left] at (0,0.2) {$n$};
  \node[above right,font=\footnotesize] at (3.7,\m) {$a_{mm}$}; \node[above right,font=\footnotesize] at (6.3,\m) {$a_{mj}$};
  \node[above right,font=\footnotesize] at (3.7,\ii) {$a_{im}$}; \node[above right,font=\footnotesize] at (6.3,\ii) {$a_{ij}$};
  \fill (3.7,\m) circle (0.05); \fill (6.3,\m) circle (0.05); \fill (3.7,\ii) circle (0.05); \fill (6.3,\ii) circle (0.05);
  % right-hand side vector
  \draw (\n+0.8,0) rectangle (\n+1.8,\n);
  \node[font=\footnotesize] at (\n+1.3,\m) {$b_m$}; \node[font=\footnotesize] at (\n+1.3,\ii) {$b_i$};
  \draw[->] (\n+1.8,\m) -- (\n+2.6,\m) -- (\n+2.6,\ii) -- (\n+1.85,\ii);
  \node[right,font=\footnotesize] at (\n+2.6,{(\m+\ii)/2}) {$\times\Bigl(-\dfrac{a_{im}}{a_{mm}}\Bigr)$};
\end{tikzpicture}

\end{fig}

\topic{Algorithm for backward substitution: $\mat{U}\mat{x}=\mat{c}\ \rightarrow\ \mat{x}$}
\begin{lstlisting}[language=Fortran,numbers=none,basicstyle=\small\ttfamily\itshape,mathescape=true]
$b_n$= $b_n/a_{nn}$
do i=n-1, 1, -1
    $b_i$=($b_i$- $\sum_{j=i+1}^{n} a_{ij}b_j$ )/$a_{ii}$
end do
\end{lstlisting}

If $\mat{A}$ is symmetric and assume that only its upper triangular coefficients $a_{ij}$ ($j\geq i$) are stored, then the above algorithm for upper triangularization needs to be modified to be of the form:
\begin{lstlisting}[language=Fortran,numbers=none,basicstyle=\small\ttfamily\itshape,mathescape=true]
do m=1, n-1
    do i=m+1, n
        r=$\boldsymbol{a_{mi}}$/$a_{mm}$
        $b_i$= $b_i$-r$b_m$
        do j=$\boldsymbol{i}$, n
            $a_{ij}$=$a_{ij}$ - r $a_{mj}$
        end do
    end do
end do
\end{lstlisting}
Actually, there are only two slight changes; the reasons for these two changes are as follows:
\begin{enumerate}
  \item $a_{im}\ \rightarrow\ a_{mi}$ because $\mat{A}$ is symmetric and its coefficients are available for the upper triangular part.
  \item \texttt{do j=m,n} $\rightarrow$ \texttt{do j=i,n} because $a_{ij} = a_{ji}$, ($j = m, i-1$) when $\mat{A}$ is symmetric even during the upper triangularization process; you may verify this by carrying out an example.
\end{enumerate}
Note that in the process of upper triangularization, if $a_{mm} = 0$, then the above algorithms will fail obviously. In this case, we need to interchange the $m$-th row and a certain row after it, say the $i$-th row, where $a_{im}\neq 0$; there are several ways to do it and this process is called \emph{pivoting}.

\section{Some widely used solvers}
\begin{enumerate}
  \item Banded solver, which is discussed in class and used in our program.
  \item Skyline solver, which will be introduced in the course of Finite Element Methods next semester.
  \item Frontal solver (see Sec.~11.5 of textbook).
\end{enumerate}

\section{Condition number}
$\operatorname{Cond}(\mat{A}) =$ condition number of $\mat{A} = \|\mat{A}^{-1}\|\,\|\mat{A}\|$; when $\mat{A}$ is symmetric such as the structural stiffness matrix $\Kmat$ and $\|\ \| = L_2$ norm, then it can be shown that
\begin{equation}
  \operatorname{cond}(\Kmat) = \frac{\max|\lambda_i|}{\min|\lambda_i|}
\end{equation}
where $\lambda_i$ are eigenvalues of $\Kmat$.

The significance of $\operatorname{cond}(\Kmat)$ lies in the fact that it is an indicator for estimating the relative error of $\bs{\Delta}$ from those of $\Kmat$ and $\Pvec$. Specifically,
\begin{equation}
  \frac{\|\delta\bs{\Delta}\|}{\|\bs{\Delta}\|} \leq
  \frac{\operatorname{cond}(\Kmat)}{1-\operatorname{cond}(\Kmat)\dfrac{\|\delta\Kmat\|}{\|\Kmat\|}}
  \left(\frac{\|\delta\Kmat\|}{\|\Kmat\|} + \frac{\|\delta\Pvec\|}{\|\Pvec\|}\right)
\end{equation}
If $\operatorname{cond}(\Kmat)\|\delta\Kmat\|/\|\Kmat\|$ is much less than 1, the above equation can be simplified to
\begin{equation}
  \frac{\|\delta\bs{\Delta}\|}{\|\bs{\Delta}\|} \leq \operatorname{cond}(\Kmat)\left(\frac{\|\delta\Kmat\|}{\|\Kmat\|} + \frac{\|\delta\Pvec\|}{\|\Pvec\|}\right)
\end{equation}
Notice that $\operatorname{cond}(\Kmat)\geq 1$. When $\operatorname{cond}(\Kmat)$ is very large, say larger than 10000, the system is called \emph{ill-conditioned}. In this case, the accuracy of $\bs{\Delta}$ is usually much less than that of $\Pvec$ and $\Kmat$.

A main cause of ill-conditioning is large differences in the stiffnesses of parts of the structure. For example, a stiff structure on flexible supports can give problems in this direction. However, a flexible structure on stiff supports tends not to suffer from such difficulties. For more examples on ill-conditioned structures, refer to Macleod (1990)\footnote{Macleod, I.A. (1990), \emph{Analytical Modelling of Structural Systems}, Ellis Horwood, London.}, page 117.

\section{Decomposition methods}
If $\mat{A}$ is symmetric (all eigenvalues are real), the two commonly used methods for solving $\mat{A}\mat{x}=\mat{b}$ are the \emph{Cholesky method} (appropriate for only positive definite matrix) and the \emph{triple-factoring method} (called root-free Cholesky method in the textbook). Their algorithms are given in Sec.~\ref{sec:cholesky-triple}.

\begin{example}[Cholesky decomposition]
\begin{equation*}
  \Kmat = \bmat{4 & 6 & 10 & 4\\ 6 & 13 & 13 & 6\\ 10 & 13 & 27 & 2\\ 4 & 6 & 2 & 72}
  = \bmat{L_{11} & 0 & 0 & 0\\ L_{21} & L_{22} & 0 & 0\\ L_{31} & L_{32} & L_{33} & 0\\ L_{41} & L_{42} & L_{43} & L_{44}}
    \bmat{L_{11} & L_{21} & L_{31} & L_{41}\\ 0 & L_{22} & L_{32} & L_{42}\\ 0 & 0 & L_{33} & L_{43}\\ 0 & 0 & 0 & L_{44}} = \mat{L}\mat{L}^\tr
\end{equation*}
\begin{align*}
  4 &= L_{11}^2 \rightarrow L_{11} = 2\\
  6 &= L_{11}L_{21} \rightarrow L_{21} = 3\\
  10 &= L_{11}L_{31} \rightarrow L_{31} = 5\\
  4 &= L_{11}L_{41} \rightarrow L_{41} = 2\\
  13 &= L_{21}^2 + L_{22}^2 \rightarrow L_{22} = 2\\
  13 &= L_{21}L_{31} + L_{22}L_{32} \rightarrow L_{32} = -1\\
  6 &= L_{21}L_{41} + L_{22}L_{42} \rightarrow L_{42} = 0\\
  27 &= L_{31}^2 + L_{32}^2 + L_{33}^2 \rightarrow L_{33} = 1\\
  2 &= L_{31}L_{41} + L_{32}L_{42} + L_{33}L_{43} \rightarrow L_{43} = -8\\
  72 &= L_{41}^2 + L_{42}^2 + L_{43}^2 + L_{44}^2 \rightarrow L_{44} = 2
\end{align*}
Therefore, $\mat{L} = \bmat{2 & 0 & 0 & 0\\ 3 & 2 & 0 & 0\\ 5 & -1 & 1 & 0\\ 2 & 0 & -8 & 2}$.
\end{example}

\begin{example}[Triple-factoring decomposition]
\begin{equation*}
  \Kmat = \bmat{1 & 2 & 3\\ 2 & 2 & 2\\ 3 & 2 & 5} = \mat{L}_1\mat{D}\mat{L}_1^\tr
\end{equation*}
\begin{equation*}
  \mat{D}\mat{L}_1^\tr = \bmat{D_1 & 0 & 0\\ 0 & D_2 & 0\\ 0 & 0 & D_3}\bmat{1 & L_{21} & L_{31}\\ 0 & 1 & L_{32}\\ 0 & 0 & 1}
  = \bmat{D_1 & D_1L_{21} & D_1L_{31}\\ 0 & D_2 & D_2L_{32}\\ 0 & 0 & D_3}
\end{equation*}
\begin{equation*}
  \Kmat = \bmat{1 & 2 & 3\\ 2 & 2 & 2\\ 3 & 2 & 5} = \mat{L}_1(\mat{D}\mat{L}_1^\tr)
  = \bmat{1 & 0 & 0\\ L_{21} & 1 & 0\\ L_{31} & L_{32} & 1}\bmat{D_1 & D_1L_{21} & D_1L_{31}\\ 0 & D_2 & D_2L_{32}\\ 0 & 0 & D_3}
\end{equation*}
\begin{align*}
  1 &= D_1\\
  2 &= D_1L_{21} \rightarrow L_{21} = 2\\
  3 &= D_1L_{31} \rightarrow L_{31} = 3\\
  2 &= D_1L_{21}^2 + D_2 \rightarrow D_2 = -2\\
  2 &= L_{21}D_1L_{31} + D_2L_{32} \rightarrow L_{32} = 2\\
  5 &= D_1L_{31}^2 + D_2L_{32}^2 + D_3 \rightarrow D_3 = 4
\end{align*}
Therefore $\mat{L}_1 = \bmat{1 & 0 & 0\\ 2 & 1 & 0\\ 3 & 2 & 1}$ and $\mat{D} = \bmat{1 & 0 & 0\\ 0 & -2 & 0\\ 0 & 0 & 4}$; $\Kmat$ is indefinite.
\end{example}

\section{Cholesky and triple-factoring methods}
\label{sec:cholesky-triple}
\textbf{Subroutines}: \texttt{DECOM} and \texttt{SOLVER}; solve $\Kmat\mat{U} = \Pvec$ for $\mat{U}$.

\topic{Factorization of $\Kmat$ (\texttt{DECOM}):}
\begin{equation}
  \Kmat = \mat{L}\mat{D}\mat{L}^\tr
\end{equation}
where $\mat{L}$ is a lower triangular matrix, and $\mat{D}$ a diagonal matrix; i.e.,
\begin{equation}
  \mat{L} = \bmat{L_{11} & & & \\ L_{21} & L_{22} & & \\ \vdots & & & \\ L_{n1} & L_{n2} & \cdots & L_{nn}}
  \quad\text{and}\quad
  \mat{D} = \bmat{D_1 & & & \\ & D_2 & & \\ & & \ddots & \\ & & & D_n}
\end{equation}
\begin{equation}
  \Kmat\mat{U} = \mat{L}\underline{\mat{D}\mat{L}^\tr\mat{U}} = \mat{L}\underline{\mat{c}} = \Pvec
\end{equation}

\topic{Two-step solution process (\texttt{SOLVER}):}
\begin{enumerate}
  \item First solve $\mat{L}\mat{c} = \Pvec$ for $\mat{c}$ by forward substitution, i.e.
    \begin{align}
      c_1 &= \frac{1}{L_{11}}P_1\\
      c_i &= \frac{1}{L_{ii}}\Bigl(P_i - \sum_{j=1}^{i-1}L_{ij}c_j\Bigr)\quad(i = 2,\cdots,n)
    \end{align}
  \item Then, solve $\mat{D}\mat{L}^\tr\mat{U} = \mat{c}$ or equivalently $\mat{L}^\tr\mat{U} = \mat{D}^{-1}\mat{c}$ for $\mat{U}$ by backward substitution, i.e.
    \begin{align}
      U_n &= \frac{1}{L_{nn}}\Bigl(\frac{c_n}{D_n}\Bigr)\\
      U_i &= \frac{1}{L_{ii}}\Bigl(\frac{c_i}{D_i} - \sum_{j=i+1}^{n}L_{ji}U_j\Bigr)\quad(i = n-1, n-2,\cdots,1)
    \end{align}
\end{enumerate}

\topic{Methods for obtaining $\mat{L}$ and $\mat{D}$:}
\begin{itemize}
  \item Unknowns: $\frac{n(n+1)}{2}$ entries of $\mat{L}$ and $n$ entries of $\mat{D}$.
  \item Equations: $\frac{n(n+1)}{2}$ equations from $\Kmat = \mat{L}\mat{D}\mat{L}^\tr$ (note that $\Kmat$ is symmetric).
\end{itemize}
We therefore have at least $n$ free choices in assigning the values of $\mat{L}$ and $\mat{D}$. We introduce two commonly used methods herein.

\begin{enumerate}[leftmargin=*]
  \item \textbf{Cholesky method}\\
    Set $D_i = 1$ ($i = 1,2,\cdots,n$); we can then determine $\mat{L}$ column by column using the formula:
    \begin{align}
      L_{jj} &= \Bigl(K_{jj} - \sum_{l=1}^{j-1}L_{jl}L_{jl}\Bigr)^{\frac12}\\
      L_{ij} &= \frac{1}{L_{jj}}\Bigl(K_{ij} - \sum_{l=1}^{j-1}L_{il}L_{jl}\Bigr)\quad(i = j+1, j+2,\cdots,n)
    \end{align}
    Note that for $j = 1$ set the summation terms to zero.

    Drawbacks:
    \begin{enumerate}[label=(\alph*)]
      \item Works efficiently for \textbf{positive definite} matrix only, because otherwise some of the $L_{ij}$ will be imaginary.
      \item Square root requires more computational efforts.
    \end{enumerate}

  \item \textbf{Triple-factoring method}\\
    Set $L_{ii} = 1$ ($i = 1,2,\cdots,n$); we can determine $D_j$ and then $L_{ij}$ columnwise by the formula:
    \begin{align}
      D_j &= K_{jj} - \sum_{l=1}^{j-1}L_{jl}D_lL_{jl}\\
      L_{ij} &= \frac{1}{D_j}\Bigl(K_{ij} - \sum_{l=1}^{j-1}L_{il}D_lL_{jl}\Bigr)\quad(i = j+1, j+2,\cdots,n)
    \end{align}
    Note that:
    \begin{enumerate}[label=(\alph*)]
      \item Works even for $\Kmat$ that is not positive definite.
      \item Determinant of $\Kmat$: $\operatorname{Det}(\Kmat) = \operatorname{Det}(\mat{L}\mat{D}\mat{L}) = D_1D_2\cdots D_n$.
      \item Operation counts including multiplication and division: $O(n^3/6)$ if $\Kmat$ is full, which is the same as that of the Cholesky method. For a banded $\Kmat$ with $\text{MHB} = m$, the operation counts are of $O(m^2n/2)$.
    \end{enumerate}
\end{enumerate}
